% Source: https://artofproblemsolving.com/community/c3h3167360_2023_lmt_spring_guts_round_p1p15_lexington_mathematical_tournament
% Stable identity: aops-topic-3167360-q04
If the representation has length 2, then $143=d(b+1)$ with $1\le d<b$. Since 143=11\cdot 13, the valid pairs are $(d,b)=(1,142),(11,12)$; the divisor 13 would give base 10 and an invalid digit. Length 4 has value $d(b+1)(b^2+1)$. Bases 2,3,4 yield factors 15,40,85, none dividing 143; base 5 already exceeds 143. Length 6 allows only base 2, whose six-one value 63 does not divide 143; base 3 is already too large. Length at least 8 has minimum 255. The sum is $142+12=\boxed{154}$.
